\subsubsection*{Division mit Rest, Euklidischer Algorithmus}


Unter welcher Bedingung ist die Gleichung $a \cdot x + b \cdot y = d$ lösbar?

\begin{lemma}[Lemma von Bezout]
Gegeben $a, b \in \Z$. Dann existieren Zahlen $x, y \in \Z$ mit
\[ a \cdot x + b \cdot y = \ggT(a, b) \qquad (*) \]
\end{lemma}

\begin{korollar}[Folgerung]
Die Gleichung $a \cdot x + b \cdot y = d$ ist genau dann lösbar, wenn $\ggT(a, b) \mid d$.
\end{korollar}

\subsubsection*{Hauptsatz der Arithmetik, Primfaktorzerlegung}

Alle Zahlen aus $\boldsymbol{N}$.

\begin{satz}[Hauptsatz der Arithmetik]
Jede natürliche Zahl $a > 1$ ist eindeutig (bis auf die Reihenfolge) ein Produkt von Primzahlpotenzen:
\[ a = p_1^{n1} \cdot p_2^{n2} \cdot \ldots \cdot p_k^{nk} \]
\end{satz}

Es ist schwer, für eine große Zahl $a$ diese Zerlegung zu gewinnen.\\
(Sieb des Eratosthenes)

\begin{definition}[ggT, kgV]
Sei $R(a, b) = \{ r \mid r \mid a \text{ und } r \mid b \}$ die Menge der gemeinsamen Teiler von $a, b$. Dann ist
\[ \boldsymbol{\ggT(a, b) = \max \{ r \mid r \in R(a, b) \}} \]
äquivalent dazu ist
\[ r \in R(a, b) \quad \Rightarrow \quad r \mid \ggT(a, b) \]
Zwei Zahlen $a, b$ mit $\ggT(a, b) = 1$ heißen teilerfremd.
\[ \boldsymbol{\kgV(a, b) = \min \{ s \mid a \mid s \text{ und } b \mid s \}} \]
\[ \kgV(a, b) \cdot \ggT(a, b) = a \cdot b \]
\end{definition}

Zur Bestimmung des $\ggT(a, b)$ wird die Primfaktorzerlegung nicht benötigt.

\begin{satz}[Division mit Rest]
Sei $a > b > 0$ gegeben. Dann sind die Zahlen $q, r$ mit
\[ a = q \cdot b + r \qquad 0 \leq r < b \]
eindeutig bestimmt.
\end{satz}

\begin{algorithmus}[Euklidischer Algorithmus]
Der \textbf{Euklidische Algorithmus} berechnet $\ggT(a, b)$.
\end{algorithmus}

\begin{beispiel}
$a = 1965$, $\;b = 225$
\begin{align*}
1965 &= 8 \cdot 225 + 165 \\
225 &= 1 \cdot 165 + 60 \\
165 &= 2 \cdot 60 + 45 \\
60 &= 1 \cdot 45 + 15 \\
45 &= 3 \cdot 15 + 0
\end{align*}
$\ggT(a, b)$ ist der letzte von Null verschiedene Zahl in der Folge der Reste, also
\[ \ggT(1965, 225) = 15. \]
\end{beispiel}

Allgemein:
\begin{align*}
a &= q_0 \cdot b + r_1 \\
b &= q_1 \cdot r_1 + r_2 \\
&\ \ \ldots \\
r_{i-1} &= q_i \cdot r_i + r_{i+1} = (q_i \ \ 1) \begin{pmatrix} r_i \\ r_{i+1} \end{pmatrix} \\
&\ \ \ldots \\
r_{n-2} &= q_{n-1} \cdot r_{n-1} + r_n, \qquad r_n > 0 \\
r_{n-1} &= q_n \cdot r_n + 0
\end{align*}

\begin{satz}[Euklid]
\[ \ggT(a, b) = r_n \]
\end{satz}

\begin{proof}
denn: $r_n \mid a$, $\;r_n \mid b$ und jeder Teiler von $a$ und $b$ teilt $r_n$.
\end{proof}

\begin{proof}[Beweis des Lemmas von Bezout]
Sei $Q_i = \begin{pmatrix} q_i & 1 \\ 1 & 0 \end{pmatrix}$. Mit $Q_i$ schreibt man
\[ \begin{pmatrix} r_{i-1} \\ r_i \end{pmatrix} = Q_i \begin{pmatrix} r_i \\ r_{i+1} \end{pmatrix} = Q_i\, Q_{i+1} \begin{pmatrix} r_{i+1} \\ r_{i+2} \end{pmatrix}, \]
also mit
\[ Q = Q_0 \ldots Q_n = \begin{pmatrix} s & u \\ t & v \end{pmatrix} \]
\[ \Rightarrow \quad \begin{pmatrix} a \\ b \end{pmatrix} = \begin{pmatrix} q_0 & 1 \\ 1 & 0 \end{pmatrix} \ldots \begin{pmatrix} q_n & 1 \\ 1 & 0 \end{pmatrix} \begin{pmatrix} r_n \\ 0 \end{pmatrix} = Q \begin{pmatrix} r_n \\ 0 \end{pmatrix} = \begin{pmatrix} s \\ t \end{pmatrix} r_n \]
Es ist $\operatorname{Det} Q_i = -1$, also
\[ \operatorname{Det} Q = s\,v - t\,u = (-1)^{n+1} \]
Im wichtigsten Fall, $a, b$ teilerfremd, also $r_n = 1$, entnimmt man eine Lösung von $(*)$ $\operatorname{Det} Q$. Andernfalls muss $\operatorname{Det} Q$ noch mit $r_n$ multipliziert werden.
\[ a \cdot v - b \cdot u = s\, r_n \cdot v - t\, r_n \cdot u = \pm r_n \]
Damit ist das Lemma von Bezout bewiesen.
\end{proof}

\begin{beispiel}[Fortsetzung]
$a = 1965$, $\;b = 225$,\\
Gesucht: Lösung der Gleichung $1965\, x + 225 \cdot y = 15$
\begin{align*}
Q &= \begin{pmatrix} 8 & 1 \\ 1 & 0 \end{pmatrix} \begin{pmatrix} 1 & 1 \\ 1 & 0 \end{pmatrix} \begin{pmatrix} 2 & 1 \\ 1 & 0 \end{pmatrix} \begin{pmatrix} 1 & 1 \\ 1 & 0 \end{pmatrix} \begin{pmatrix} 3 & 1 \\ 1 & 0 \end{pmatrix} \\
&= \begin{pmatrix} 131 & 35 \\ 15 & 4 \end{pmatrix} = \begin{pmatrix} s & u \\ t & v \end{pmatrix}
\end{align*}
\[ \operatorname{Det} Q = 131 \cdot 4 - 15 \cdot 35 = -1 \quad | \cdot (-15) \]
\[ \Rightarrow \quad 1965 \cdot (-4) + 225 \cdot 35 = 15 \]
\end{beispiel}

\subsubsection*{Wie schnell ist der EA?}

Für alle $1 \leq i \leq n - 1$ gilt die Abschätzung
\[ r_{i-1} = q_i \cdot r_i + r_{i+1} \geq 1 \cdot r_i + r_{i+1} > 2\, r_{i+1} \]
\begin{align*}
&\text{Insbesondere} && b > 2\, r_2 > 2^2 r_4 > 2^3 r_6 > \ldots > 2^n\, r_{2n} \\
&\text{Aus} && 2^n > b \quad \text{folgt} \quad r_{2n} = 0 \\
&\text{m.a.W.} && n > \log_2 b \quad \text{folgt} \quad r_{2n} = 0
\end{align*}

\begin{satz}
Der EA benötigt höchstens $O(\log_2 b)$ Iterationen bis zum $\ggT(a, b)$.
\end{satz}

$\log_2 b$ ist die Anzahl der Binärstellen von $b$.\\
Der EA hat die Komplexität $O(\log b)$ bzw.\ linear in der Anzahl der Dezimal- (Binär-) stellen von $b$.

\subsubsection*{Homomorphismen von Gruppen}

\begin{definition}
Seien $(G,+)$, $(H,+)$ zwei Gruppen und $f : G \to H$ eine Abbildung von $G$ nach $H$. Falls
\[ f(g + g') = f(g) + f(g') \qquad \text{für alle } g, g' \in G, \]
dann heißt $f$ ein (\textbf{Gruppen-})\textbf{Homomorphismus} von $G$ nach $H$.
\end{definition}

Die Struktur von $G$ findet sich also in $H$ wieder, wenn auch meist etwas gröber.

\begin{beispiel}
$G = \Z$, $\;H = \Zm{m}$, $\;f(g) = [g]$
\end{beispiel}

\begin{definition}
Falls $G, H$ Ringe sind und zusätzlich
\[ f(g \cdot g') = f(g) \cdot f(g') \qquad \text{für alle } g, g' \in G, \]
gilt, ist $f$ ein \textbf{Ringhomomorphismus}.

Ist $f$ zusätzlich bijektiv, so spricht man von einem \textbf{Isomorphismus}.

Gruppen (Ringe) $G, H$ heißen \textbf{isomorph}, falls ein Isomorphismus $f : G \to H$ existiert. Man schreibt dann $G \cong H$.
\end{definition}

\subsubsection*{Der Ring $\Zm{m} \times \Zm{n}$}

Auf dem kartesischen Produkt $\Zm{m} \times \Zm{n}$ sind Addition und Multiplikation wie folgt erklärt
\begin{align*}
(a_1, a_2) + (b_1, b_2) &= (a_1 + b_1, a_2 + b_2) \\
(a_1, a_2) \cdot (b_1, b_2) &= (a_1 \cdot b_1, a_2 \cdot b_2)
\end{align*}
$\Zm{m} \times \Zm{n}$ ist ein Ring mit dem Einselement $(1, 1)$.
\[ f : \Zm{m \cdot n} \to \Zm{m} \times \Zm{n}, \quad \text{mit} \]
\[ f(z) = (z \bmod m,\ z \bmod n) \]
ist ein Homomorphismus.

\begin{beispiel}
Die Ringe $\Zm{12}$, $\;\Zm{2} \times \Zm{6}$ sind nicht isomorph, denn
\[ f(0) = f(6) = (0,0) \qquad f(1) = f(7) = (1,1) \]
$f$ ist nicht surjektiv, da die Bilder $(0, 1)$, $(1, 0)$ nicht vorkommen;\\
$f$ ist nicht injektiv, da die Bilder $(0, 0)$, $(1, 1)$ je zwei Urbilder haben.
\end{beispiel}

Der Chinesische Restsatz gibt die Bedingung an unter der die Ringe
\[ \Zm{m \cdot n}, \quad \Zm{m} \times \Zm{n} \]
isomorph sind.
